% Policy Content Hiding section body.
\subsection{Policy Content Hiding (PCH)}
\label{sec:eval-pch}

\noindent\textbf{Property.} PCH requires that the public-chain observer cannot recover the \emph{content} of a policy---its DUO data-use terms or the ICD-10 disease codes it references---better than chance.

\noindent\textbf{Scope of the labels.} The attack targets the ten DUO data-use terms of the corpus---general research use (GRU), disease-specific research (DS), research-specific restrictions (RS), population origins or ancestry research (POA), geographical restriction (GS), clinical care use (CC), non-commercial use (NCU), health or medical or biomedical research (HMB), no methods development (NMDS), and genetic studies only (GSO)---and disease codes aggregated to the ten ICD-10 chapters. This label restriction, its statistical-power rationale, and why category-level granularity is a conservative choice are detailed in Section~\ref{sec:eval-dataset}.

\noindent\textbf{Supervised attack.} Each DUO term defines a one-vs-rest binary task, on which we run the multi-seed pipeline of Section~\ref{sec:eval-method}.
Table~\ref{tab:pch-terms} reports the per-term verdicts: every DUO task stays below the $0.55$ AUC line, with mean AUC between $0.477$ and $0.540$; significant-seed fractions run up to $23\%$ (GRU), short of the $50\%$ gate, and with the AUC leg unmet none approaches the failure criterion.
The same pipeline applied to the ten ICD-10 chapters (one-vs-rest over the ICD-bearing rows only, median test $n{=}127$ per chapter; Table~\ref{tab:pch-icd10}) passes nine chapters and fails one.
Chapters ch1, ch2, and ch4 sit just above the AUC line (mean AUC $0.558$--$0.570$) but reproduce in at most $23\%$ of seeds, so they are boundary cases rather than stable channels.
Chapter ch5 (diseases of the nervous system) meets both legs: mean AUC $0.667$ with $53\%$ of seeds significant.
We report it as a FAIL of the criterion and examine below what the classifier separates.

\noindent\textbf{What the ch5 classifier separates.} The two gas features are the only non-constant observables, and the access-gas measurement is itself the sum of two parts: the target that the TPS draws uniformly at random from the padding band for that execution (Section~\ref{sec:design}) and the gas that the padded call consumes beyond it, which on Amoy is a near-constant $4{,}661$--$4{,}750$ gas for all $1773$ executions.
The observer sees only the sum, but the capture log records every target, so we can ask how much of the separability survives once the target is subtracted.
None does.
Over all $20$ PCH tasks the rank-based AUC of the access gas coincides with the AUC of the target alone to three decimals, and the AUC of the residual (gas minus target) lies between $0.47$ and $0.52$; for ch5 the gas and target AUCs are both $0.601$ and the residual AUC is $0.472$.
The $34$ ch5 policies simply drew targets whose mean lies $33{,}000$ gas above that of the other $388$ ICD-bearing policies, a $3.2\sigma$ deviation for a uniform draw of that size; because each coverage-layer policy is executed once, the $30$ seeds resample the same $34$ fixed draws, which is why the deviation reproduces across seeds.
Rerunning the of-record protocol unchanged on the residual in place of the access gas gives mean AUC $0.556$ with $10\%$ significant seeds (logistic regression) and $0.561$ with $17\%$ (random forest), a PASS; rerunning it on the target alone reproduces the FAIL ($0.686$, $63\%$).
The padded gas therefore carries a policy-independent random draw that happened to correlate with one chapter in this corpus, not information about the chapter; Section~\ref{sec:eval-method} quantifies how often the criterion produces such a false positive.
